\section*{الفصل \ref{ch:b3:adaptive-immunity} --- المناعة التكيّفية والتلقيح}

\begin{solution}{exo:b3:adaptive-immunity:1}
\omterm{def:b3:adaptive-immunity:lymphocytes}{لمفاوية} واحدة، ونوعية واحدة (نوسال وليدربرغ: كل خلية أفرزت
ضدًا \omterm{def:b3:adaptive-immunity:lymphocytes}{لمستضد} واحد فقط)؛ والمستقبل يُصنع قبل لقاء \omterm{def:b3:adaptive-immunity:lymphocytes}{المستضد}
(تونيغاوا: المورّثة يُعاد ترتيبها في \omterm{def:b3:adaptive-immunity:lymphocytes}{اللمفاوية البائية}، لا
استجابةً \omterm{def:b3:adaptive-immunity:lymphocytes}{للمستضد})؛ \omterm{def:b3:adaptive-immunity:lymphocytes}{والمستضد} ينتخب النسائل الملائمة ويوسّعها
إلى خلايا منفِّذة وخلايا ذاكرة؛ والنسائل المتفاعلة مع الذات
تُحذف أثناء النماء.
\end{solution}

\begin{solution}{exo:b3:adaptive-immunity:2}
سلسلتان ثقيلتان وسلسلتان خفيفتان، في كل منها نطاق متغير في
الطرف ونطاقات ثابتة وراءه؛ وموقعا ربط متطابقان تكوّنهما ست
عُرى شديدة التغير؛ وساق Fc تقرؤها خلايا أخرى وتقرؤها
\omterm{def:b3:innate-immunity:complement}{المتمّمة}. وIgM: أول الأضداد، خماسي، يفعّل \omterm{def:b3:innate-immunity:complement}{المتمّمة}. وIgG:
\omterm{def:b3:adaptive-immunity:antibody}{ضد} الدم والنسج الرئيس، \omterm{def:b3:innate-immunity:complement}{للأوبسنة} والتعديل، ويعبر المشيمة.
وIgA: ثنائي يُفرز على الأسطح المخاطية وفي الحليب. وIgE:
مرتبط بالخلايا البدينة، \omterm{def:b3:adaptive-immunity:antibody}{ضد} الطفيليات، ووسيط الأرجية. وIgD:
على اللمفاويات البائية الساذجة، وظيفته المعروفة ثانوية.
\end{solution}

\begin{solution}{exo:b3:adaptive-immunity:3}
الصنف~I: كل الخلايا ذات النواة؛ وببتيدات من بروتينات
العصارة الخلوية يقطعها البروتيازوم؛ من 8--10 بقايا؛ تقرؤها
اللمفاويات التائية السامة CD8. والصنف~II: الخلايا التغصنية،
والبلعميات، واللمفاويات البائية؛ وببتيدات من بروتينات
التُقطت وهُضمت في جسيمات داخلية؛ من 13--25 بقية؛ تقرؤها
اللمفاويات التائية المساعدة CD4.
\end{solution}

\begin{solution}{exo:b3:adaptive-immunity:4}
$R_{0}$ هو متوسط عدد الحالات التي تنتجها حالة واحدة في
جمهرة قابلة كلها؛ والعتبة $p_{c} = 1 - 1/R_{0}$. فالنكاف
\qty{80}{\%}؛ والسعال الديكي \qty{92}{\%}؛ وإنفلونزا 1918 \qty{60}{\%}.
\end{solution}

\begin{solution}{exo:b3:adaptive-immunity:5}
$50\times 20\times 5 = 5000$ سلسلة ثقيلة، و$60\times 4 = 240$ خفيفة:
أي $1.2\times 10^{6}$ توفيقة؛ ومع \omterm{def:b3:adaptive-immunity:antibody}{التنوع الوصلي} $1.2\times
10^{10}$ --- أي عشرة أضعاف عدد اللمفاويات البائية، فيعبّر
الحيوان عن عُشر ما يقدر على صنعه على الأكثر، وتحمل كل خلية
تقريبًا مستقبلًا فريدًا.
\end{solution}

\begin{solution}{exo:b3:adaptive-immunity:6}
$1000\times 10^{-3} = 1$ طفرة في الخلية في الانقسام؛ و$20$ في
عشرين انقسامًا. وبين $10^{5}$ خلية يقع $2\times 10^{6}$ حدث
طفر على $3000$ بديل أحادي ممكن: فكل منها يُعاين مئات المرات،
مع التوافيق الثنائية والثلاثية فوق ذلك --- \omterm{def:b3:adaptive-immunity:germinal-centre}{فالمركز الإنتاشي}
يستكشف جوار المستقبل استكشافًا مستوفى.
\end{solution}

\begin{solution}{exo:b3:adaptive-immunity:7}
تُنتخب اللمفاويات التائية في \omterm{def:b3:adaptive-immunity:tolerance}{التوتة} لتتعرف الببتيدات على
أليلات MHC الخاصة بالمضيف؛ وأليلات شخص آخر لها أخاديد مختلفة
الشكل تحمل ببتيدات مختلفة باتجاهات مختلفة، فلا يُرى الببتيد
الفيروسي نفسه. وجزيئات MHC الغريبة في طعم تتعرفها هي نفسها
نسبة كبيرة من اللمفاويات التائية على أنها ``MHC ذاتي بببتيد
غريب''، فيكون الرفض. ويدوم تعدد الأشكال لأن جمهرة ذات
أليلات كثيرة تعرض عينة أوسع من ببتيدات أي ممرض ولا يقدر
ممرض على الإفلات من الجميع؛ والمتغايرو الزيجة يعرضون ضعف
المدى فيُحابون.
\end{solution}

\begin{solution}{exo:b3:adaptive-immunity:8}
(أ) لا \omterm{def:b3:adaptive-immunity:antibody}{إعادة تركيب V(D)J}: فلا لمفاويات بائية ولا تائية، أي
عوز مناعي مشترك شديد. (ب) لا AIRE: فتخفق خلايا \omterm{def:b3:adaptive-immunity:tolerance}{التوتة} في
عرض مستضدات النسج، وتفلت اللمفاويات التائية المتفاعلة مع
الذات، فتنشأ \omterm{def:b3:adaptive-immunity:tolerance}{مناعة ذاتية} \omterm{def:b3:adaptive-immunity:antibody}{ضد} أعضاء كثيرة. (ج) لا \omterm{def:b3:adaptive-immunity:tolerance}{FoxP3}: فلا
لمفاويات تائية نظيمية، وتنشأ \omterm{def:b3:adaptive-immunity:tolerance}{مناعة ذاتية} وأرجية قاتلة باكرة
البدء. (د) لا AID: فلا تبديل صنف ولا طفور مفرط --- أي IgM
وافر، ولا يكاد يوجد IgG ولا IgA، وأضداد منخفضة الألفة،
وأخماج متكررة. (ه) لا لمفاويات تائية CD4: فلا مساعدة
للمفاويات البائية ولا للبلعميات، واستجابات ضدّية رديئة
وقابلية للأخماج الانتهازية --- وهو عوز الإيدز المناعي.
\end{solution}

\begin{solution}{exo:b3:adaptive-immunity:9}
$p_{c} = 1 - 1/6 = 0.83$؛ والتغطية $0.83/0.9 = 93\,\%$. وعند
تغطية \qty{80}{\%} تكون نسبة المناعيين $0.72$، و$R_{\text{eff}} = 6\times 0.28
= 1.7$. وبين القابلين ينتشر الوباء كأن $R_{0}$ هو
$1.7$ (إذ لا يفعل المناعيون إلا تخفيف التماسّات)، وتعطي
علاقة الحجم النهائي $s_{\infty} - \ln s_{\infty}/1.7 = 1$ قيمة $s_{\infty}
\approx 0.31$: أي نحو \qty{69}{\%} من القابلين، و\qty{19}{\%} من
الجمهرة، يُخمجون.
\end{solution}

\begin{solution}{exo:b3:adaptive-immunity:10}
من $\mathrm{d}s/\mathrm{d}t = -\beta si$ و$\mathrm{d}i/\mathrm{d}t
= \beta si - \gamma i$ نجد $\mathrm{d}i/\mathrm{d}s = -1 + \gamma/(\beta s)$،
فيكون $i + s - (\ln s)/R_{0}$ محفوظًا؛ ومع $s = 1$ و$i = 0$ في
البداية و$i = 0$ و$s = s_{\infty}$ في النهاية يكون $s_{\infty} - \ln
s_{\infty}/R_{0} = 1$. والحلول: $R_{0} = 1.5$، و$s_{\infty} = 0.42$
(\qty{58}{\%} مخموجين)؛ و$2$، و$0.20$ (\qty{80}{\%})؛ و$3$، و$0.06$
(\qty{94}{\%})؛ و$5$، و$0.007$ (\qty{99.3}{\%}). وليس الجميع: إذ
يهبط عدد التكاثر الفعال دون الواحد مع نفاد القابلين بينما
يبقى بعضهم، فينطفئ الوباء قبل أن يبلغهم.
\end{solution}

\begin{solution}{exo:b3:adaptive-immunity:11}
$1.5^{n} = 1000$: أي $n = \ln 1000/\ln 1.5 \approx 17$ تحسينًا
متتاليًا. ومع طفرة واحدة في الخلية في الانقسام وواحدة في
مئة تحسّن، يجب فحص مئة خلية على الأقل في الجولة لإيجاد
واحدة؛ وعمليًا آلاف، إذ تهدم نصف الطفرات الربط ويكون
الانتخاب مشوَّشًا. وسبع عشرة جولة من انقسامين لكل واحدة
تستغرق نحو تسعة أيام --- وهي الأسبوعان إلى الثلاثة
المرصودان، بحسبان أوجه القصور. والدورات تفصل الطفر (المنطقة
الداكنة) عن الاختبار (المنطقة الفاتحة) حتى يُحكم على كل
متغير في مواجهة \omterm{def:b3:adaptive-immunity:lymphocytes}{المستضد} قبل أن يُسمح له بالتكاثر، تمامًا
كما تتناوب الخوارزمية الجينية التباين والانتخاب.
\end{solution}

\begin{solution}{exo:b3:adaptive-immunity:12}
النساء: عدة مورّثات مناعية تقع على الصبغي X، وبعضها (TLR7)
يفلت من التعطيل ويُعبَّر عنه من النسختين، والإستروجينات
تعزز بقاء اللمفاويات البائية --- والاختبار: الرجال الذين
لديهم X إضافي (كلاينفلتر) يكون خطر الذئبة لديهم أنثويًا،
وهو ما يقع فعلًا. \omterm{def:b3:adaptive-immunity:mhc}{وHLA}: الأليلات المرتبطة تعرض ببتيدات
الذات ذات الصلة عرضًا جيدًا، أو تخفق في عرضها في \omterm{def:b3:adaptive-immunity:tolerance}{التوتة}
فلا تُحذف النسائل --- والاختبار: الفئران المحوَّرة بالأليل
البشري يصيبها المرض حين تُعطى \omterm{def:b3:adaptive-immunity:lymphocytes}{المستضد}. وأما الارتفاع: فقلة
التعرض الميكروبي الباكر وتبدل الميكروبيومات (الصادات،
والحمية، والولادة القيصرية) تعطي لمفاويات تائية نظيمية أقل
والتهابًا أكثر --- والاختبارات: الفئران المؤهبة للسكري إذا
رُبّيت خالية من الجراثيم أو على الصادات أصابها السكري أكثر،
وأبناء المهاجرين يكتسبون حدوث البلد الجديد في جيل واحد.
\end{solution}

\begin{solution}{pb:b3:adaptive-immunity:1}
\textbf{1.} $6000\times 320 = 1.9\times 10^{6}$؛ و$\times 3\times 10^{4}
= 5.8\times 10^{10}$.
\textbf{2.} $10^{11}$ خلية في مقابل $5.8\times 10^{10}$ ممكنة:
فقد يسع الجسم كل مستقبل مرتين تقريبًا، غير أن النسائل غير
متساوية ومعظم المتواليات لا يقع صنعها قط، فيعبّر كل جسم عن
عينة كبيرة لكن ناقصة --- ولا يتشارك شخصان إلا نسبة صغيرة من
مستقبلاتهما.
\textbf{3.} $10^{11}/10^{5} = 10^{6}$ خلية؛ وفي عقدة فيها $10^{8}$،
نحو $1000$.
\textbf{4.} $10^{9}/10^{5} = 10^{4}$ سليفة --- أي أقل بمئة ضعف؛
وللوليد أيضًا جريبات وخلايا تغصنية غير ناضجة، ولا ذاكرة له،
ومساعدة تائية محدودة، فتكون استجاباته بطيئة صغيرة، ويسدّ
IgG الأمومي الفجوة.
\textbf{5.} إنها تمتد على وصلتي V--D وD--J، حيث تُزال
النوكليوتيدات وتُضاف عشوائيًا: فمتواليتها ليست مرمَّزة في
أي مكان من الجينوم. وإذ تقع بين السلسلتين في مركز موقع
الربط، تكون العروة الأكثر ملامسةً للحاتمة.
\textbf{6.} قبل \omterm{def:b3:adaptive-immunity:lymphocytes}{المستضد} يجب أن يكون التنوع واسعًا حتى يربط
شيء ما كل ما يأتي؛ وبعد \omterm{def:b3:adaptive-immunity:lymphocytes}{المستضد} يُنفَق الجهد خيرًا في تنويع
المستقبلات القليلة المعروف أنها تربط. فبحث خشن ثم بحث دقيق
أنجع بكثير من أي منهما وحده.
\textbf{7.} سبعة أيام هي $21$ انقسامًا: أي $100\times 2^{21} = 2.1\times
10^{8}$ خلية؛ و\qty{30}{\%} منها، أي $6.3\times 10^{7}$ \omterm{def:b3:adaptive-immunity:response}{خلية بلازمية}.
\textbf{8.} $6.3\times 10^{7}\times 2000\times \num{86400} = 1.1\times
10^{16}$ جزيء في اليوم؛ و$\times 1.5\times 10^{5}\times 1.66\times
10^{-24}$ غرام $= \qty{2.7}{mg}$؛ وفي \qty{3}{L}، \qty{0.9}{mg/L} في اليوم.
\textbf{9.} \qty{27}{mg} في \qty{3}{L}: أي \qty{9}{mg/L} $= \qty{0.009}{g/L}$،
أي جزء من ألف من مجموع IgG --- \omterm{def:b3:adaptive-immunity:antibody}{فالضد} لأي \omterm{def:b3:adaptive-immunity:lymphocytes}{مستضد} واحد نسبة
صغيرة من الكل.
\textbf{10.} المئة ضعف هي $\log_{2}100 = 6.6$ عمر نصف، أي نحو
$140$ يومًا. وسنوات \omterm{def:b3:adaptive-immunity:antibody}{الضد} تأتي من الخلايا البلازمية الطويلة
العمر في النقي، التي تفرز باستمرار، ومن خلايا الذاكرة التي
تُستثار من جديد عند إعادة التعرض.
\textbf{11.} تسعة انقسامات: $10^{5}\times 512 = 5\times 10^{7}$ خلية في
اليوم 3 --- أي ألف ضعف $5\times 10^{4}$ للأولية في اليوم 3،
وربع ما بلغته الأولية في اليوم 7 فقط.
\textbf{12.} تقع VDJ المعاد ترتيبها بجوار C$_{\mu}$، فيكون أول
نسخة IgM؛ والتبديل إلى C$_{\gamma}$ يحتاج إلى مساعدة تائية
وإلى AID، وهما يستغرقان أيامًا. وخلايا الذاكرة قد بدّلت
فعلًا، فتصنع ذريتها IgG فورًا.
\textbf{13.} $1000\times 10^{-3} = 1$ في الانقسام؛ و$20$ على مدى عشرين.
\textbf{14.} واحدة في مئة؛ و$10^{4}\times 0.01 = 100$ خلية محسَّنة
في كل انقسام.
\textbf{15.} $1.5^{n} = 1000$: أي $n \approx 17$؛ و$17\times \qty{12}{h}
\approx 8.5$ يوم.
\textbf{16.} نصف البنات تفقد الربط: فمن دون انتخاب، لا يبقى
رابطًا بعد عشرين انقسامًا إلا $0.5^{20}
\approx 10^{-6}$ من السلالة. فلا بد أن تزيل المنطقة الفاتحة
الخاسرات في كل دورة حتى لا تمضي إلا المستقبلات العاملة
والمحسَّنة.
\textbf{17.} تترك الجرعة الأولى خلايا ذاكرة بائية مرتفعة
الألفة مبدَّلة الصنف؛ وتزرع الجرعة الثانية بها مراكز
إنتاشية جديدة، فيستأنف الطفور والانتخاب من نقطة بداية أعلى،
فتكون الأضداد الخارجة أفضل بعد.
\textbf{18.} بعض المخموجين يصنعون في النهاية أضدادًا \omterm{def:b3:adaptive-immunity:antibody}{ضد} مواقع
محفوظة لا يقدر \omterm{def:b3:virology:virus}{الفيروس} على تغييرها (أضداد معادِلة واسعة
الطيف)، بعد سنين من النضج. \omterm{def:b3:adaptive-immunity:germinal-centre}{والمركز الإنتاشي} يمكن توجيهه:
فمتوالية من المُمنِّعات تُشرك أولًا سلائف السلالة الجرثومية
الصحيحة ثم تكافئ الربط بالموقع المحفوظ قد تقود النضج على
ذلك المسار في أشهر بدل سنين.
\textbf{19.} $p_{c} = 1 - 1/15 = 93.3\,\%$. وبجرعة واحدة: $0.933/0.93 =
100.4\,\%$ --- وهو غير قابل للبلوغ؛ وبجرعتين: $0.933/0.97 = 96.2\,\%$.
\textbf{20.} المناعيون $0.90\times 0.97 = 0.873$؛ و$R_{\text{eff}} = 15\times
0.127 = 1.9$: فالحصبة تدور.
\textbf{21.} $60\,000\times (0.10 + 0.90\times 0.03) = 7600$ قابل في كل
فوج ولادات؛ وبعد خمس سنين $38\,000$ --- أي مجمع يجد فيه
وافد مخموج من القابلين في تماسّاته ما يكفي لإدامة السلاسل.
\textbf{22.} بين القابلين، تعطي $s_{\infty} - \ln s_{\infty}/1.9 =
1$ قيمة $s_{\infty} \approx 0.23$: أي نحو \qty{77}{\%} منهم، ونحو
$29\,000$ حالة.
\textbf{23.} المناعيون $0.95\times 0.97 = 0.92$؛ و$R_{\text{eff}} = 15\times
0.08 = 1.2$ --- ولا يزال فوق الواحد؛ فالعتبة تحتاج إلى
\qty{96}{\%} بجرعتين، والحصبة هي المرض الذي يعاقب على آخر بضع
في المئة.
\textbf{24.} لا يحمي الرُّضّعَ إلا مناعة من حولهم (\omterm{def:b3:adaptive-immunity:antibody}{والضد}
الأمومي بضعة أشهر): فإذا لم تتكون سلاسل الانتقال لم تبلغهم
حالة. وعند تغطية \qty{85}{\%} تكون نسبة المناعيين $0.82$
و$R_{\text{eff}} = 2.6$: فتجري الفاشيات، ويُخمج الرُّضّع، وهم
غير مُلقَّحين وأشد عرضة.
\textbf{25.} الذخيرة نحو $6\times 10^{10}$؛ و\qty{2.7}{mg} من IgG في
اليوم من خلايا استجابة واحدة البلازمية؛ وتغطية \qty{96}{\%}
بجرعتين.
\end{solution}
